% Created 2026-01-05 Mon 22:46
% Intended LaTeX compiler: pdflatex
\documentclass[11pt]{article}
\usepackage[utf8]{inputenc}
\usepackage[T1]{fontenc}
\usepackage{graphicx}
\usepackage{longtable}
\usepackage{wrapfig}
\usepackage{rotating}
\usepackage[normalem]{ulem}
\usepackage{amsmath}
\usepackage{amssymb}
\usepackage{capt-of}
\usepackage{hyperref}
\author{B. Orçun Özkablan}
\date{\today}
\title{}
\hypersetup{
 pdfauthor={B. Orçun Özkablan},
 pdftitle={},
 pdfkeywords={},
 pdfsubject={},
 pdfcreator={Emacs 30.2 (Org mode 9.7.11)}, 
 pdflang={English}}
\begin{document}

\tableofcontents

\section{Haskell Cheatsheet}
\label{sec:orgd5fb38c}

\subsection{🧱 Basics}
\label{sec:org897e40f}
\begin{center}
\begin{tabular}{ll}
Syntax & Description\\
\hline
x = 5 & Define a constant named `x`\\
-- comment & Single-line comment\\
let x = 5 in x + 1 & Define local variable in expression\\
where & Local definitions in a function body\\
\end{tabular}
\end{center}
\subsection{🧮 Functions}
\label{sec:orge90e753}
\begin{center}
\begin{tabular}{ll}
Syntax & Description\\
\hline
f x = x + 1 & Define a function\\
add a b = a + b & Function with two arguments\\
\x -> x + 1 & Lambda (anonymous function)\\
f :: Int -> Int & Type annotation\\
f x = g (h x) & Function composition using nesting\\
f = g . h & Function composition using (.)\\
\end{tabular}
\end{center}
\subsection{📊 Types}
\label{sec:orgd5739a9}
\begin{center}
\begin{tabular}{ll}
Syntax & Description\\
\hline
Int & Fixed-size integer\\
Integer & Arbitrary-precision integer\\
Float & Floating-point number\\
Double & Double-precision float\\
Bool & Boolean (`True`, `False`)\\
Char & Character literal (e.g., 'a')\\
String & List of characters\\
{[}a] & List of elements of type `a`\\
(a, b) & Tuple with two elements\\
\end{tabular}
\end{center}
\subsection{📚 Pattern Matching}
\label{sec:orgecf484c}
\begin{center}
\begin{tabular}{ll}
Syntax & Description\\
\hline
f [] = 0 & Match an empty list\\
f (x:xs) = x & Match head and tail of a list\\
f (a, b) = a + b & Destructure tuple\\
\end{tabular}
\end{center}
\subsection{🧩 Tuples \& Lists}
\label{sec:org2987e9f}
\begin{center}
\begin{tabular}{ll}
Syntax & Description\\
\hline
(x, y) & Tuple\\
fst (x, y) & First element of tuple\\
snd (x, y) & Second element of tuple\\
{[}1, 2, 3] & List\\
{[}1..5] & Range from 1 to 5\\
head xs & First element\\
tail xs & All elements except first\\
length xs & Length of list\\
x:xs & Prepend `x` to list `xs`\\
xs ++ ys & Concatenate two lists\\
\end{tabular}
\end{center}
\subsection{🔁 Control \& Recursion}
\label{sec:orgdcf046c}
\begin{center}
\begin{tabular}{ll}
Syntax & Description\\
\hline
if cond then a else b & Conditional expression\\
case x of \ldots{} & Pattern match expression\\
factorial 0 = 1 & Base case for recursion\\
factorial n = n * factorial(n-1) & Recursive definition\\
\end{tabular}
\end{center}
\subsection{🧠 Guards}
\label{sec:orged9989c}
\begin{center}
\begin{tabular}{ll}
Syntax & Description\\
\hline
f x & x > 0 = \ldots{} & Conditional logic in functions\\
 & otherwise = \ldots{} & Default (else) case\\
\end{tabular}
\end{center}
\subsection{🧪 Maybe Type}
\label{sec:org4a776a8}
\begin{center}
\begin{tabular}{ll}
Syntax & Description\\
\hline
Maybe a & Optional value\\
Just 5 & A value is present\\
Nothing & No value present\\
case x of Just n -> \ldots{} & Pattern match on Maybe\\
\end{tabular}
\end{center}
\subsection{🎛 Higher-Order Functions}
\label{sec:org5762965}
\begin{center}
\begin{tabular}{ll}
Syntax & Description\\
\hline
map f xs & Apply function to each element\\
filter p xs & Keep elements that satisfy predicate\\
foldl (+) 0 xs & Left fold (reduce from left)\\
zipWith (+) [1,2] [3,4] & Element-wise operation on two lists\\
\end{tabular}
\end{center}
\subsection{🔍 GHCi Commands}
\label{sec:org6fe3816}
\begin{center}
\begin{tabular}{ll}
Command & Description\\
\hline
ghci & Start interactive Haskell shell\\
:l file.hs & Load a file into GHCi\\
:r & Reload the file\\
:t expr & Show the type of an expression\\
:i name & Show info about a function/type\\
\end{tabular}
\end{center}
\end{document}
